\section{2025 Mathematical Physics}

\begin{exercise}
\label{physics:2025:1}
Bead on a rotating hoop. Analyze equilibria and bifurcation.
\end{exercise}

\begin{solution}
\textbf{Prerequisites:}
\begin{itemize}
    \item \textbf{Lagrangian in Rotating Frame:} $L = T - V$.
    \item \textbf{Effective Potential:} $V_{\text{eff}} = - \frac{1}{2}m\omega^2 R^2 \sin^2\theta - mgR\cos\theta$.
\end{itemize}

\textbf{Steps:}

1. \textbf{Equation of Motion:} $\ddot{\theta} = \sin\theta(\omega^2\cos\theta - g/R)$. 
2. \textbf{Critical Frequency:} For $\omega < \Omega = \sqrt{g/R}$, the only stable equilibrium is at the bottom ($\theta=0$).
3. \textbf{Bifurcation:} For $\omega > \Omega$, the bottom becomes unstable ($V_{\text{eff}}'' < 0$), and a new set of stable equilibria appears at $\cos\theta_0 = g/R\omega^2$.
4. \textbf{Oscillations:} Near the stable point $\theta_0$, small vibrations have frequency $\omega \sin\theta_0$.

\textbf{Intuition:}
Imagine sitting in a spinning laundry machine. If it spins slowly, you stay at the bottom due to gravity. If it spins fast enough, the "centrifugal push" overcomes gravity, and you get flung to the side. This sudden switch between "staying down" and "climbing up" is a classic example of symmetry breaking.
\end{solution}

\begin{exercise}
\label{physics:2025:2}
Wave packet in $z$-direction. Longitudinal fields.
\end{exercise}

\begin{solution}
\textbf{Prerequisites:}
\begin{itemize}
    \item \textbf{Gauss's Law:} $\nabla \cdot \mathbf{E} = 0$.
    \item \textbf{Circular Polarization:} $\mathbf{E}_\perp \propto (\hat{x} \pm i\hat{y})$.
\end{itemize}

\textbf{Steps:}

1. \textbf{Constraint:} $\partial_x E_x + \partial_y E_y + \partial_z E_z = 0$.
2. \textbf{Emergence:} If the beam has a finite width (e.g. Gaussian profile), $\partial_x E_x \neq 0$. To satisfy Gauss's law, there MUST be a longitudinal field $E_z \approx \frac{i}{k} (\nabla \cdot \mathbf{E}_\perp)$.
3. \textbf{Significance:} This small $E_z$ field is necessary to give the photon its expected spin and orbital angular momentum along the beam axis.

\textbf{Intuition:}
Light is usually described as being purely transverse. However, if you "squeeze" a beam to a finite size, the field lines have to curve back. This curving creates a tiny component of the electric field pointing in the direction the light is traveling. You can't have a localized beam without this longitudinal "swirl."
\end{solution}

\begin{exercise}
\label{physics:2025:3}
Shifted harmonic oscillator $V(x) = \frac{1}{2}m\omega^2 x^2 - q\mathcal{E}x$.
\end{exercise}

\begin{solution}
\textbf{Prerequisites:}
\begin{itemize}
    \item \textbf{Completing the Square:} Transforming potential energy forms.
    \item \textbf{Displacement Operator:} $D(\alpha) = e^{\alpha a^\dagger - \alpha^* a}$.
\end{itemize}

\textbf{Steps:}

1. \textbf{Shift:} $V(x) = \frac{1}{2}m\omega^2 (x - x_0)^2 + \Delta E$, where $x_0 = q\mathcal{E}/m\omega^2$.
2. \textbf{Eigenvalues:} The energies are $E_n = \hbar\omega(n+1/2) - q^2\mathcal{E}^2/2m\omega^2$. The field just lowers all levels uniformly.
3. \textbf{Overlap:} The new ground state is a shifted version of the old one. The probability of staying in the old ground state is $|\langle 0 | 0' \rangle|^2 = \exp(-x_0^2/2l^2)$.

\textbf{Intuition:}
Applying an electric field to an oscillator is like pulling the end of the spring. The spring doesn't get "softer" or "harder" (frequency stays the same), it just has a new natural resting point. In the quantum world, this shift means you are "partially" in many of the old energy levels at once.
\end{solution}

\begin{exercise}
\label{physics:2025:4}
Ising model on a triangle. Frustration and thermodynamics.
\end{exercise}

\begin{solution}
\textbf{Prerequisites:}
\begin{itemize}
    \item \textbf{Partition Function:} $Z = \sum e^{-\beta E}$.
    \item \textbf{Frustration:} Inability to satisfy all bond interactions simultaneously.
\end{itemize}

\textbf{Steps:}

1. \textbf{Enumeration:} $2^3 = 8$ states. $E = -J \sum \sigma_i \sigma_j$.
2. \textbf{Z Calculation:} $Z = 2e^{3\beta J} \cosh(3\beta h) + 6e^{-\beta J} \cosh(\beta h)$.
3. \textbf{Heat Capacity:} From $U = -\partial_\beta \ln Z$, $C_v$ shows a peak but no phase transition (since it is a finite system).
4. \textbf{Antiferromagnetism:} If $J < 0$, the system cannot make all pairs happy, leading to a degenerate "frustrated" ground state.

\textbf{Intuition:}
Imagine three people who all want to disagree with each other. Person A disagrees with B, B disagrees with C... but now C is forced to either agree with A or B! This "social deadlock" in a triangle of interactions is why frustrated magnets behave so strangely compared to regular ones.
\end{solution}

\begin{exercise}
\label{physics:2025:5}
Gravitational waves from binary BH merger.
\end{exercise}

\begin{solution}
\textbf{Prerequisites:}
\begin{itemize}
    \item \textbf{Quadrupole Formula:} $h \sim \ddot{Q}/r$.
    \item \textbf{Kepler's Law:} $r \Omega^2 \sim 1/r^2$.
\end{itemize}

\textbf{Steps:}

1. \textbf{Moment:} $Q_{jk} = \sum m x_j x_k$. For circular motion, $Q \propto \cos(2\Omega t)$.
2. \textbf{Power radiated:} $P \propto \dddot{Q}^2 \propto \Omega^6 r^4$. Using Kepler's law, $P \propto 1/r^5$.
3. \textbf{Chirp:} As the system loses energy, $r$ shrinks, forcing $\Omega$ to increase to stay in orbit. This leads to the characteristic "chirp" frequency rise.

\textbf{Intuition:}
A binary black hole is like a massive mixer stirring the fabric of space itself. The ripples it sends out carry away energy. To compensate for the loss, the black holes must fall closer together, which makes them whirl even faster, sending out waves even more violently until they finally collide.
\end{solution}

\begin{exercise}
\label{physics:2025:6}
Conformal scalar field with curvature coupling $\xi R \phi^2$.
\end{exercise}

\begin{solution}
\textbf{Prerequisites:}
\begin{itemize}
    \item \textbf{Conformal Transformation:} $g \to \Omega^2 g$.
    \item \textbf{Ricci Scalar Change:} $\bar{R} = R/\Omega^2 + \dots$
\end{itemize}

\textbf{Steps:}

1. \textbf{Invariance:} Kinetic term $\partial \phi^2$ isn't conformally invariant on its own in curved space.
2. \textbf{Coupling:} Setting $\xi = (d-2)/4(d-1)$ (which is $1/6$ in 4D) exactly cancels the curvature terms generated during the transformation.
3. \textbf{Effective Mass:} In de Sitter space ($R=12H^2$), the coupling $\xi R \phi^2$ looks like a mass term $m^2 = \xi R = 2H^2$.

\textbf{Intuition:}
Usually, being "massless" means you travel at the speed of light and your physics is the same regardless of scale. In a curved universe, the stretching of space normally breaks this. By adding a specific interaction between the field and the curvature, we "re-balance" the scales so the field remains truly scale-less.
\end{solution}
